\documentclass[11pt]{article}
\usepackage[a4paper, margin=1in]{geometry}
\usepackage{amsmath, amssymb}
\usepackage{enumitem}
\usepackage{fancyhdr}
\usepackage{xcolor}

\pagestyle{fancy}
\fancyhf{}
\lhead{CSE 400: Fundamentals of Probability in Computing}
\rhead{Lecture Scribe}
\cfoot{\thepage}

\newcommand{\solution}{
	\vspace{0.5cm}
	\noindent\textbf{\textcolor{blue}{Solution:}}
	\vspace{0.2cm}
}

\title{
	\normalsize School of Engineering and Applied Science (SEAS), Ahmedabad University \\
	\vspace{0.2cm}
	\textbf{CSE 400: Fundamentals of Probability in Computing}\\
	\Large Lecture Scribe Submission
}

\author{}
\date{}

\begin{document}
\maketitle

\vspace{-2cm}
\begin{center}
	\begin{tabular}{ll}
		\textbf{Name:} & {Kush Kelaiya \hspace{2.5in}} \\ [1.5ex]
		\textbf{Enrollment No:} & {AU2440018 \hspace{2.5in}} \\ [1.5ex]
		\textbf{Email:} & {kush.k2@ahduni.edu.in \hspace{2.1in}} \\ [1.5ex]
		\textbf{Group: } & s1\_g1\_fin \\ [1.5ex]
		\textbf{Lecture:} & 6 \\ [1.5ex]
		\textbf{Date of Submission:} & {February 7\textsuperscript{th}, 2026 (11:59PM) \hspace{1.2in}}
	\end{tabular}
\end{center}

\hrule
\vspace{0.5cm}

\section{Random Variables}
\subsection{Motivation and Concept}
A random variable (RV) is a function that maps outcomes of a random experiment (sample space) to real numbers.

\begin{itemize}
	\item The \textbf{distribution of a random variable} can be visualized as a bar diagram.
	\item The \textbf{x-axis} represents the values that the random variable can take.
	\item The \textbf{height of the bar at a value $a_i$} is the probability of $Pr\left[ X=a_i \right]$.
	\item Each probability is computed from the probability of the corresponding event in the sample space.
\end{itemize}

\section{Types of Random Variable}
\subsection{Discrete Random Variable}
A random variable is \textbf{discrete} if:
\begin{itemize}
	\item It has \textbf{countable support}.
	\item It is described using a \textbf{Probability Mass Function (PMF)}.
	\item Probabilities are assigned to \textbf{single values}.
	\item Each possible value has \textbf{strictly positive probability}.
\end{itemize}
\subsection{Continuous Random Variable}
A random variable is continuous if:
\begin{itemize}
	\item It has \textbf{uncountable support}.
	\item It is described using a \textbf{Probability Density Function (PDF)}.
	\item Probabilities are assigned to \textbf{intervals of values}.
	\item Each individual value has \textbf{zero probability}.
\end{itemize}

\section{Examples of Random Variables}
\subsection*{Example 1: Tossing Three Fair Coins}
\textbf{Experiment:} Toss 3 fair coins.\\
Let $Y$ denote the \textbf{number of heads} that appear.
\begin{itemize}
	\item Possible values of $Y$ : $\left\{ 0, 1, 2, 3 \right\}$
\end{itemize}
Probability Calculations
\begin{itemize}
	\item $Pr(Y=0) = Pr(t, t, t) = \frac{1}{8}$
	\item $Pr(Y=1) = Pr(t, t, h) + Pr(t, h, t) + Pr(h, t, t) = \frac{3}{8}$
	\item $Pr(Y=2) = Pr(t, h, h) + Pr(h, t, h) + Pr(h, h, t) = \frac{3}{8}$
	\item $Pr(Y=3) = Pr(h, h, h) = \frac{1}{8}$
\end{itemize}
Thus,
$$Pr(Y=0) = \frac{1}{8},\,Pr(Y=1) = \frac{3}{8},\,Pr(Y=2) = \frac{3}{8},\,Pr(Y=3) = \frac{1}{8}$$
Since $Y$ must take one of these values:
$$1 = Pr\left( \bigcup_{i=0}^3(Y=i) \right) = \sum_{i=0}^{3}Pr(Y=i)$$

\section{Probability Mass Function (PMF)}
\subsection{Definition}
A random variable that can take on \textbf{at most a countable number of possible values} is said to be \textbf{discrete}.\\\\
Let $X$ be a discrete random variable with range (possible values)
$$R_x = \left\{ x_1, x_2, x_3, \dots \right\}$$

(where the set is \textbf{finite or countably infinity}).\\
The function
$$P_X(X_K) - Pr(X=\mathsf{X}_K), k = 1, 2, 3, \dots$$
is called the \textbf{Probability Mass Function (PMF)} of $X$.
\subsection{PMF Property}
Since $X$ must take \textbf{one of the values} $X_K$:
$$\sum_{k=1}^{\infty}P_X(X_K)=1$$

\section{PMF - Worked Example}
\subsection*{Example}
The probability mass function of a random variable $X$ is given by:
$$p(i) = c\dfrac{\lambda^i}{i!}, i = 0, 1, 2, \dots$$
where $\lambda > 0$.\\
\textbf{Step 1: Find the constant C}\\
Using the PMF condition:
\begin{align*}
	\sum_{i=0}^{\infty}p(i) &= 1\\\\
	c\sum_{i=0}^{\infty}\dfrac{\lambda^i}{i!} &= 1\\
\end{align*}
Since,
\begin{align*}
	\epsilon^\lambda = \sum_{i=0}^{\infty}\dfrac{\lambda^i}{i!}
\end{align*}
we get:
$$c\epsilon^\lambda = 1 \to c = \epsilon^{-\lambda}$$
\textbf{Step 2:} Find Pr[X=0]
	$$Pr[X=0] = c\dfrac{\lambda^0}{0!} = \epsilon^{-\lambda}$$
\textbf{Step 3:} Find Pr[X>2]\\
\begin{align*}
Pr[X>2] = 1 - Pr[X\leq2]\\
= 1-1 {Pr[X=0] + Pr[X=1] + Pr[X=2]}\\
= 1 - \left( \epsilon^{-\lambda} + \lambda\epsilon^{-\lambda} + \dfrac{\lambda^2}{2}\epsilon^{-\lambda} \right)
\end{align*}

\section{Bayes' Theorem}
\subsection{Recap}
Using conditional probability:
$$Pr(A\cap B_i) = Pr(B\,|\,A)Pr(A)$$
We obtain:
$$
Pr(B_i\,|\,A) = \dfrac{Pr(A\,|\,B_i)Pr(B_i)}{\sum_{n}^{j=1}pr(A\,|\,B_j)Pr(B_j)}
$$
This is known as \textbf{Bayes' Formula (Proposition 3.1)}.
\subsection{Terminology}
\begin{itemize}
	\item $Pr(B_i)$: \textbf{A priori probability}\\
	(formed from presupposed or self-evident models)
	\item $Pr(B_i\,|\,A)$: \textbf{Posterior probability}
	(computed after observing event $A$)
\end{itemize}

\section{Bayes' Theorem - Example 1}
\subsection*{Auditorium with 30 Rows}
\begin{itemize}
	\item The auditorium has 30 rows.
	\item Row 1 has 11 seats, Row 2 has 12 seats, Row 3 has 13 seats, \dots, Row 30 has 40 seats.
	\item A door prize is given by
	\begin{enumerate}
		\item Randomly selecting a row (each row equally likely.)
		\item Randomly select a seat within that row (each seat equally likely).
	\end{enumerate}
\end{itemize}
\textbf{Tasks}
\begin{enumerate}
	\item Compute the probability that Seat 15 was selected given that Row 20 was selected.
	\item Compute the probability that Row 20 was selected given that Seat 15 was selected. 
\end{enumerate}
(Computation is carried out using Bayes' formula as shown in the lecture.)
\section{Bayes' Theorem - Example 2}
\subsection*{Communication System (Receiver - RX)}
\begin{itemize}
	\item A communication system sends \textbf{binary data}: 0 or 1.
	\item The receiver (RX) may make errors:
	\item \begin{itemize}
		\item A 0 may be detected as 1.
		\item A 1 may be detected as 0.
	\end{itemize}
	\item The system is described using a given \textbf{set of conditional probabilities.}
\end{itemize}
(The posterior probabilities are computed using Bayes' theorem based on the provided conditional probabilities.)

\vfill
\bigskip
\hrule
\vspace{0.2cm}
\begin{center}
\small \textit{End of Submission}
\end{center}

\end{document}